\documentclass[../../main.tex]{subfiles}
\begin{document}

\begin{center}
  {\fontsize{48pt}{58pt}\selectfont \textbf{1997}\normalsize\textbf{年度}}
\end{center}

\begin{table}[htb]
  \centering
  \caption{出題テーマと難易度}
  \begin{tabular}{|c|p{9cm}|c|}
    \hline
    \textbf{問題番号} & \textbf{テーマ} & \textbf{難易度} \\
    \hline
    第1問 & 不等式の条件（楕円と直線） & \\
    \hline
    第2問 & 数列の極限（区分求積とはさみうち） & \\
    \hline
    第3問 & 整数問題（単位分数の分解） & \\
    \hline
    第4問 & 三角形の面積公式と最小値 & \\
    \hline
  \end{tabular}
\end{table}

\begin{table}[htb]
  \centering
  \caption{解答}
  \renewcommand{\arraystretch}{1.4}
  \begin{tabular}{|c|c|p{9cm}|}
    \hline
    \textbf{問題番号} & \textbf{小問} & \textbf{答え} \\
    \hline
    第1問 & - & $a+b\ge\sqrt{2}$（$a=0$なら$b\ge1$，$b=0$なら$a\ge1$，$a=b=0$も可） \\
    \hline
    第2問 & (1) & $\log 2$ \\
    \cline{2-3}
    & (2) & 証明問題 \\
    \hline
    第3問 & (1) & $(x,y)=(3,6),(4,4),(6,3)$ \\
    \cline{2-3}
    & (2) & 証明問題 \\
    \hline
    第4問 & (1) & 証明問題 \\
    \cline{2-3}
    & (2) & 最小値 $\dfrac{3\sqrt{3}}{28}$ \\
    \hline
  \end{tabular}
\end{table}

\paragraph{試験形態}
4問．試験時間は要確認．

\end{document}
